\documentclass[10pt,a4paper]{article}
\usepackage[top=0.95in,bottom=0.95in,left=1.1in,right=1.1in]{geometry}
\usepackage{amsmath,amssymb}
\usepackage[colorlinks=true,linkcolor=black,citecolor=black,urlcolor=blue]{hyperref}
\usepackage[protrusion=true,expansion=false]{microtype}
\usepackage[T1]{fontenc}
\setlength{\parskip}{3pt}
\setlength{\parindent}{0pt}

\title{\textbf{Is the Quantum Circuit Necessary?}\\
\large A Classifiability-Predicted Ablation of Quantum-Weight-Generation for Physics-Informed Neural Networks}
\author{mesh.applied}
\date{July 2026}

\begin{document}
\maketitle

\begin{abstract}
Quantum-Train is an architecture in which a parameterized quantum circuit (PQC) is trained to generate the weights of a classical neural network, after which the quantum circuit is discarded and only the classical weights are deployed. Its authors prove that a classical algorithm cannot reproduce this generation process efficiently in general, unless the polynomial hierarchy collapses, but they also state that this proof does not hold for circuits with one-dimensional or geometrically local structure, where a classical tensor network already performs the same task efficiently. No published work tests which side of this boundary a Quantum-Train circuit falls on when used for a physics-informed neural network (PINN) solving a partial differential equation (PDE), a task class where weight generation has never been evaluated. We propose to close this gap for the two-dimensional Burgers' equation: classify the specific circuit used for weight generation against existing classical-simulability criteria, state the resulting prediction before training, then run a matched-parameter-count ablation between the quantum generator and two classical generators under identical training conditions. The result is a direct test of whether the necessity proof's stated exception applies here, and a reusable procedure for checking the same claim in any future quantum-weight-generation paper.
\end{abstract}

\section{The Problem}

For a physics-informed neural network solving a partial differential equation, does a Quantum-Train-style weight generator produce measurably different training or generalization behavior than a classical generator with the same number of trainable parameters?

Quantum-Train works as follows \cite{liu2024}. A parameterized quantum circuit with $N$ qubits and rotation angles $\theta$ is trained, then measured in the computational basis. The measurement outcomes are passed through a small classical linear layer to produce the weight matrix $W_c$ of a target classical network:
\[
q_k = \left|\langle k|\,U(\theta)\,|0\rangle^{\otimes N}\right|^2 \quad (k=0,\dots,2^N-1), \qquad W_c = Aq + b.
\]
After training, $\theta$ is fixed, the circuit is run once to extract $W_c$, and the quantum circuit is discarded; only $W_c$ is deployed. In one reported case this compresses a classical network of $M=6690$ parameters into 728 quantum circuit parameters, a ratio of 10.8 percent \cite{liu2024}.

Liu et al.\ prove that, unless SampBQP equals SampBPP (an equality that would collapse the polynomial hierarchy, a foundational assumption in computational complexity theory), no classical algorithm can efficiently reproduce this weight-generation process in general \cite{liu2024}. This is a theorem, not an empirical claim. The same paper states the theorem's boundary directly: matrix product states, a classical tensor-network representation, already simulate the relevant quantum states efficiently when the circuit is restricted to one dimension or a geometrically local Hamiltonian, so the necessity argument does not apply to that circuit class \cite{liu2024}.

Whether a given weight-generating circuit falls inside or outside this boundary determines whether using it is justified by the theorem or is an unexamined assumption. Nobody has checked which side of the boundary current applications fall on for a physics-informed network. Existing Quantum-Train extensions cover image classification, reinforcement learning, flood prediction, federated learning, and typhoon trajectory forecasting \cite{liu2025typhoon}, but not a PDE-constrained network, where the loss function, training dynamics, and required parameter structure differ from classification or time-series tasks. Anyone building a Quantum-Train-based PDE solver is currently unable to state whether the quantum component of their architecture does anything a classical generator could not.

\section{Prior Work}

Liu et al.\ \cite{liu2024} introduce Quantum-Train and prove the necessity result above. Their evaluation is limited to image classification; the one-dimensional and geometrically local exception is stated but never checked against the circuits they use.

Cerezo et al.\ \cite{cerezo2025} show, across many circuit families used to avoid barren plateaus (a training failure mode in which gradients vanish exponentially as circuit size grows), that the same structural properties which prevent vanishing gradients often permit efficient classical simulation once some classical data has been collected from the quantum device. This establishes, in an adjacent part of quantum machine learning, that checking a specific circuit against classical-simulability criteria is a standard and productive audit, not a novel or speculative step.

Sweke et al.\ \cite{sweke2025} give matching necessary and sufficient conditions under which random Fourier features, a classical technique, can or cannot reproduce a PQC-based regression model. This supplies the general mathematical framework for asking, precisely, whether a classical method can replace a quantum one for a specific task, and shows that the answer depends on stated properties of the circuit rather than being decided case by case without theory.

Three recent papers approach the target application without closing the gap. Liu et al.\ \cite{liu2025typhoon} apply the exact Quantum-Train mechanism, including an ablation against classical compression baselines (pruning, weight sharing), but to typhoon trajectory forecasting, a time-series task, and against compression techniques rather than independent classical weight generators at matched parameter count. Zabihi et al.\ \cite{zabihi2026} benchmark a hybrid quantum-classical PINN on Burgers', Allen-Cahn, and Korteweg-de Vries equations, the same PDE family targeted here, but keep the PQC permanently embedded in the inference path rather than generating weights once and discarding it, so their result answers whether a persistent quantum layer helps, not whether quantum-generated weights differ from classically-generated ones. Linghu et al.\ \cite{linghu2026} test a photonic quantum representation on seven PDE benchmarks including Burgers', with frozen and shuffled ablation controls, but the photonic circuit is again a persistent part of the trained representation, not a weight generator that is discarded after training.

\section{Proposed Solution}

We combine two existing tools that have not previously been applied together, then test the combination empirically.

\begin{itemize}
\item \textbf{Classification.} Classify the entangling structure of the weight-generating circuit against the classical-simulability criteria of Cerezo et al.\ \cite{cerezo2025}, specifically whether its causal light cone stays shallow enough for a classical-shadow-assisted simulator to reproduce its measurement statistics. This produces a falsifiable prediction before any training: inside the classically simulable regime, no measurable advantage over a classical generator is expected; outside it, the outcome is open under both \cite{liu2024} and \cite{cerezo2025}.
\item \textbf{Matched baselines.} Build two classical generators with the same parameter count as the quantum generator: a random low-rank projection, and a matrix product state generator with the bond dimension the classification step identifies as sufficient.
\item \textbf{Ablation.} Train all three generators, quantum, low-rank, and matrix-product-state, under identical conditions: same physics-informed loss, same optimizer schedule, same collocation points, on the two-dimensional Burgers' equation. Measure PDE residual error, boundary and initial condition loss, generalization to a held-out spatiotemporal region, and stability across random seeds.
\end{itemize}

This differs from Zabihi et al.\ \cite{zabihi2026} and Linghu et al.\ \cite{linghu2026}, which test whether a persistently active quantum layer improves a PDE solver, by testing instead whether a quantum-generated, then frozen, set of weights differs from a classically generated set of the same size. It differs from Liu et al.'s typhoon application \cite{liu2025typhoon}, which compares against compression techniques on a time-series task, by comparing against independent weight generators at matched parameter count on a PDE task. No new quantum hardware, proof, or PDE solver is required. What is new is the combination: an existing classification criterion applied to an existing weight-generation architecture, tested against an existing PDE benchmark, under a matched-parameter experimental design that has not been run.

\section{Applications and Impact}

The primary application domain is scientific machine learning, specifically physics-informed neural networks used to solve partial differential equations.

If the quantum generator shows no measurable advantage, the result determines what any Quantum-Train-based PDE solver can honestly claim: a compression and regularization technique with a classical equivalent, not a demonstration of quantum necessity. This prevents a specific failure mode already documented in quantum machine learning, where architectures whose loss landscapes avoid barren plateaus are later shown to be classically simulable after publication \cite{cerezo2025}. If the generator does show an advantage, it is the first documented case of Quantum-Train compression outperforming a matched classical baseline on a differential-equation task rather than a classification or time-series task, which would justify further work on physics-informed applications specifically.

Beyond this architecture, the classification-then-ablation procedure is directly reusable. Any paper claiming that a quantum circuit generates classical parameters more effectively than a classical alternative, in reinforcement learning, federated learning, or large language model fine-tuning \cite{liu2025typhoon}, can be checked with the same two steps: classify the circuit against known classical-simulability criteria, then test the resulting prediction against a matched-parameter classical baseline. This turns an unverifiable claim about quantum necessity into a repeatable, falsifiable check.

In computational fluid dynamics and structural engineering, physics-informed networks already serve as fast surrogates for expensive numerical solvers. A weight-generation method proven, rather than assumed, to compress parameters without losing accuracy is directly usable for deploying PDE surrogates on memory-constrained hardware, independent of whether the generator involves a quantum circuit or its classical equivalent.

\begin{thebibliography}{10}
\footnotesize

\bibitem{liu2024}
C.-Y.\ Liu, E.-J.\ Kuo, C.-H.\ A.\ Lin, S.\ Chen, J.\ G.\ Young, Y.-J.\ Chang, M.-H.\ Hsieh.
\newblock Training Classical Neural Networks by Quantum Machine Learning.
\newblock arXiv:2402.16465, 2024.

\bibitem{cerezo2025}
M.\ Cerezo, M.\ Larocca, D.\ Garc\'ia-Mart\'in, N.\ L.\ Diaz, P.\ Braccia, E.\ Fontana, M.\ S.\ Rudolph, P.\ Bermejo, A.\ Ijaz, S.\ Thanasilp, E.\ R.\ Anschuetz, Z.\ Holmes.
\newblock Does provable absence of barren plateaus imply classical simulability?
\newblock Nature Communications 16, 7907, 2025. arXiv:2312.09121.

\bibitem{sweke2025}
R.\ Sweke, E.\ Recio-Armengol, S.\ Jerbi, E.\ Gil-Fuster, B.\ Fuller, J.\ Eisert, J.\ J.\ Meyer.
\newblock Potential and limitations of random Fourier features for dequantizing quantum machine learning.
\newblock Quantum 9, 1640, 2025. arXiv:2309.11647.

\bibitem{liu2025typhoon}
C.-Y.\ Liu, K.-C.\ Chen, Y.-C.\ Chen, S.\ Y.-C.\ Chen, W.-H.\ Huang, W.-J.\ Huang, Y.-J.\ Chang.
\newblock Quantum-Enhanced Parameter-Efficient Learning for Typhoon Trajectory Forecasting.
\newblock arXiv:2505.09395, 2025.

\bibitem{zabihi2026}
K.\ Zabihi, H.\ Montazeri, A.\ S.\ J.\ Suiker.
\newblock Hybrid quantum-classical physics-informed neural networks for solving nonlinear PDEs: when and where hybridization is effective?
\newblock arXiv:2606.04679, 2026.

\bibitem{linghu2026}
J.\ Linghu, H.\ Dong, Y.\ Wang.
\newblock Trainable Photonic Measurement for Physics-Informed PDE Learning.
\newblock arXiv:2606.18713, 2026.

\end{thebibliography}

\end{document}
